\section{Results}\label{sec:results}

This section reports each family's outputs as the frozen analysis wrote them, with every verdict in
the analysis's own words. Every value comes from a generated macro file or a generated table.
Section~\ref{sec:deviations} gives what changed after the freeze and how the analysis was run.
Table~\ref{tab:invalid} in Appendix~\ref{app:tables} lists every invalid window of a family cell, per
arm, with its reasons.

\subsection{Cost}\label{sec:res-cost}

Table~\ref{tab:cost} gives every cell of the cost family. Of the \CostHolm{} cells in Holm's
procedure, \CostPassed{} pass under both computations, and the two computations disagree on
\CostDisagree{} cells. C1 passes in \CountPassCOne{} of its \CountHolmCOne{} cells in Holm, C2 in
\CountPassCTwo{} of \CountHolmCTwo, and C3 in \CountPassCThree{} of \CountHolmCThree.

\begin{table}[H]
\caption{The cost family, one-port / dedicated. C1 and C2 are throughputs, so a ratio below one
favours dedicated mode; C3 is the median time to first byte, so a ratio above one does. The
intervals are the \CiPct\% BCa interval and the interval at the cell's Holm level. The verdict is the
analysis's, cut before its first parenthesis.\label{tab:cost}}
\begin{adjustwidth}{-\extralength}{0cm}
\centering\scriptsize
\begin{tabularx}{\fulllength}{@{}l l l r l l l X@{}}
\toprule
\textbf{Hyp.} & \textbf{Host, backend} & \textbf{Protocol} & \textbf{Valid} & \textbf{Median} &
\textbf{\CiPct\% interval} & \textbf{Holm interval} & \textbf{Verdict} \\
\midrule
\tabinput{../results/tables/cost}
\bottomrule
\end{tabularx}
\end{adjustwidth}
\end{table}

The two cells in Holm that fail are C3's h2c cells on L. On epoll the median ratio of the time to
first byte is \CThreeLEpollHTwoCMedian, and its \CiPct\% interval,
[\CThreeLEpollHTwoCLow, \CThreeLEpollHTwoCHigh], lies wholly above \MarginHigh. So the analysis
reports a measured cost, against one-port mode. Only \CThreeLEpollHTwoCSignXHigh{} of its
\CThreeLEpollHTwoCValid{} sessions lie below the upper margin. On io\_uring the median is
\CThreeLIouringHTwoCMedian, with the interval [\CThreeLIouringHTwoCLow, \CThreeLIouringHTwoCHigh].
Its lower end lies inside the margin, so the verdict is that equivalence is not shown, and no
measured cost is claimed there. \CThreeLIouringHTwoCSignXHigh{} of its \CThreeLIouringHTwoCValid{}
sessions lie below the upper margin. The C1 and C2 cells of h2c pass on both Linux backends. The
joint power the pilot gave h2c's three cells was \JointPowerLEpollHtwoc{} on epoll and
\JointPowerLIouringHtwoc{} on io\_uring, so these failures are not a matter of low power.

A claim of no measurable cost needs C1, C2 and C3 to pass for one protocol, backend and host. It
holds for \CostClaimCount{} configurations. They are \CostClaims. The joint power that the pilot
gave each of these claims is \LEpollHttpJointPower{} on L, except for MQTT on io\_uring with
\LIouringMqttJointPower. On W it is \WIocpHttpJointPower{} for HTTP/1.1 and \WIocpTlsJointPower{}
for TLS. As stated before the run, the TLS cells have little power to show the cost of a check of
\TlsNeed{} bytes. So their passes are not read as evidence about that check.

On W, h2c and MQTT have no claim. Their C1 cells are outside the confirmatory family, since W's
churn of these two protocols measures the generator (Section~\ref{sec:dev-before}). The pilot's
rows are consistent with that. Over the \PilotWWindows{} pilot windows of each cell, the server's CPU was busy for a
median of \PilotWHtwocServerBusy\% in h2c and \PilotWMqttServerBusy\% in MQTT. It was busy for
\PilotWHttpServerBusy\% in HTTP/1.1, whose server closes first. The generator's CPUs were
\PilotWHtwocGenCpusBusy\% and \PilotWMqttGenCpusBusy\% busy, so neither side was saturated.
Development diagnostics traced the limit to the generator's bind and close, which Windows
serializes across its threads. Their C3 cells were not resolved by the pilot. They ran at \RC{} sessions and are reported as not resolved, with no
equivalence claimed. h2c has the median \CThreeWIocpHTwoCMedian{} with the interval
[\CThreeWIocpHTwoCLow, \CThreeWIocpHTwoCHigh], and MQTT \CThreeWIocpMqttMedian{} with
[\CThreeWIocpMqttLow, \CThreeWIocpMqttHigh]. Their open-loop rates, \RateFrac{} times W's
generator-bound churn medians, were \LambdaWHtwoc{} exchanges per second for h2c and
\LambdaWMqtt{} for MQTT, against \LambdaWHttp{} for HTTP/1.1. So on W they ran at a smaller share of the server's capacity than open
HTTP/1.1. Their C2 cells pass.

No cost window on L was invalid. On W the C3 cells of HTTP/1.1, h2c and MQTT have
\CThreeWIocpHttpInvalidWindows, \CThreeWIocpHTwoCInvalidWindows{} and
\CThreeWIocpMqttInvalidWindows{} invalid windows; Table~\ref{tab:invalid} gives their arms and
reasons. The analysis clustered by lab job decides nothing. Only W's cost cells ran in more than one
job, and their clustered intervals are in Table~\ref{tab:byjob}. With two jobs such an interval
spans little more than the two jobs' medians, so it carries little information.

The secondary metrics of WL4 are in Table~\ref{tab:wl4}. In every C1 and C3 cell on L, one-port
mode's median server CPU per exchange was above dedicated mode's. In C1's HTTP/1.1 cell on epoll,
for example, it was \SecWLFourCOneLEpollHttpCpuMedian{} times, with the interval
[\SecWLFourCOneLEpollHttpCpuLow, \SecWLFourCOneLEpollHttpCpuHigh]. In C2 the medians lie on both
sides of one. The tail of the time to first byte follows the median in C3's h2c cells on L: the
ratio of \PNinetyNine{} is \SecWLFourCThreeLEpollHTwoCTtfbPNinetyNineMedian{} on epoll and
\SecWLFourCThreeLIouringHTwoCTtfbPNinetyNineMedian{} on io\_uring.

\subsection{Robustness}\label{sec:res-robust}

\runin{B1 and B2.} The hard-case runners ran every case of the frozen table, \CaseReps{} times, on
each entry of backend, detection mode and dispatch: \HardcaseLEntries{} entries on L and
\HardcaseWEntries{} on W. That is \HardcaseLRuns{} runs on L and \HardcaseWRuns{} on W, and all
\HardcasePassed{} of the \HardcaseRuns{} runs passed. They hold \HardcaseBOneDeviations{}
deviations from the frozen table and \HardcaseBTwoFailures{} failed checks of B2. Every server check
at the end of an entry was in order: \HardcaseLServerChecksOk{} of \HardcaseLServerChecks{} on L
and \HardcaseWServerChecksOk{} of \HardcaseWServerChecks{} on W. The analysis lists \BOneFailures{}
measured windows that counted a connection classified other than as its script's protocol. Its
count covers the windows that record one, the one-port windows with in-process dispatch. In the
other measured windows, a check of the server's counters per class after the analysis found only
each script's protocol and the probe. So B1 \BOneVerdict{} and B2 \BTwoVerdict. In
\HardcaseLWokeMoreThanOnce{} runs on L and \HardcaseWWokeMoreThanOnce{} on W the detection woke
more than once before it decided.

\runin{B2's distributions.} Table~\ref{tab:b2} gives the lateness of timed events and the decision
latency per backend, which the frozen text reports and does not test. They come from the same runs,
and every run of a backend that holds a value counts, whatever its case, detection mode or
dispatch. The lateness is the time from a timed event's deadline to the return of the wait that
the loop handled it after. Its median was \BTwoDistEpollLatenessMedian~ms on epoll,
\BTwoDistIouringLatenessMedian~ms on io\_uring and \BTwoDistIocpLatenessMedian~ms on IOCP. The
largest lateness, \BTwoDistEpollLatenessMax, \BTwoDistIouringLatenessMax{} and
\BTwoDistIocpLatenessMax~ms, lies below the guard band $G$ of each host, \GuardL~ms on L and
\GuardW~ms on W. The decision latency runs from the client's last write before the end of detection
to that end, on one clock. In the runs no timer ended, its median was
\BTwoDistEpollDecisionBytesMedian~$\mu$s on epoll, \BTwoDistIouringDecisionBytesMedian~$\mu$s on
io\_uring and \BTwoDistIocpDecisionBytesMedian~$\mu$s on IOCP. These runs include HC21 and HC22 with
a partial signature, whose detection ends on the client's half-close or reset, a pause after its
last write. There are \BTwoDistEpollDecisionBytesShutN{} such runs on each Linux backend and
\BTwoDistIocpDecisionBytesShutN{} on IOCP. Their latency holds that pause: at least
\BTwoDistEpollDecisionBytesShutMin~$\mu$s on epoll, \BTwoDistIouringDecisionBytesShutMin~$\mu$s on
io\_uring and \BTwoDistIocpDecisionBytesShutMin~$\mu$s on IOCP. On the Linux backends they form
the upper tail: every other run took at most \BTwoDistEpollDecisionBytesOtherMax~$\mu$s on epoll and
\BTwoDistIouringDecisionBytesOtherMax~$\mu$s on io\_uring. On IOCP \BTwoDistIocpDecisionBytesOtherAboveN{}
other runs took longer than any of them, in \BTwoDistIocpDecisionBytesOtherAboveCases, the longest
\BTwoDistIocpDecisionBytesOtherMax~$\mu$s. In the runs a timer ended, detection ends when the loop
handles the timer, so the latency holds the time from the client's last write to that handling.

\begin{table}[H]
\caption{B2's distributions per backend, over the hard-case runs: L's \BTwoDistEpollRuns{} runs on
each Linux backend and W's \BTwoDistIocpRuns{} on IOCP. The lateness is that of each timed event
past its deadline, in the runs a timer ended. The decision latency runs from the client's last write
before the end of detection to that end. A run with no such write has none: \BTwoDistEpollNoWriteN{}
per backend on L and \BTwoDistIocpNoWriteN{} on W. Its column gives \PNinetyNine, the value at
rank \NearestRankNinetyNine{} of the sorted values.\label{tab:b2}}
\centering\scriptsize
\begin{tabularx}{\textwidth}{@{}l X r r r r r@{}}
\toprule
\textbf{Backend} & \textbf{Distribution (unit)} & \textbf{Runs} & \textbf{Min} & \textbf{Median} &
\textbf{\PNinetyNineHead} & \textbf{Max} \\
\midrule
\tabinput{../results/tables/b2}
\bottomrule
\end{tabularx}
\end{table}

\runin{B3.} Table~\ref{tab:b3} gives every cell. Of the \BThreeHolm{} cells in Holm, \BThreePassed{}
pass under both computations, which disagree on \BThreeDisagree{} cells. In each pass the
competitor held at least \FootprintMargin{} times the server's counted footprint per pending
connection, by B3's test. The passes are nginx, HAProxy and Envoy in both cases, and hyper-util in
the silent case, each on epoll and on io\_uring. Each comes with its difference $D$, the
competitor's $W$ less the server's, in bytes per pending connection. The passes' medians of $D$
range from \BThreeSilentHyperUtilIouringDMedian{} bytes against hyper-util on io\_uring to
\BThreePartialHelloEnvoyEpollDMedian{} against Envoy in the partial-ClientHello case on epoll. The
server's own bounds for a pending connection are those of Section~\ref{sec:modes}. It holds no data
buffer before the first byte, the handler's buffer in replay, and at most \Bch{} of a ClientHello in
pass-through.

\begin{table}[H]
\caption{B3, competitor / server, on L. $Q$ is the ratio of counted footprints, with its median,
\CiPct\% interval and interval at Holm's level; $X$ counts the sessions with $Q$ above
\FootprintMargin. $D$ is the difference in bytes per pending connection, with its \CiPct\%
interval. ``Shown'' stands for the analysis's verdict that the competitor held at least
\FootprintMargin{} times the server's counted footprint per pending connection, by B3's
test.\label{tab:b3}}
\begin{adjustwidth}{-\extralength}{0cm}
\centering\scriptsize
\begin{tabularx}{\fulllength}{@{}l l l r l l l r l l X@{}}
\toprule
\textbf{Case} & \textbf{System} & \textbf{Backend} & \textbf{Valid} & \textbf{$Q$} &
\textbf{\CiPct\% interval} & \textbf{Holm interval} & \textbf{$X$} & \textbf{$D$} &
\textbf{\CiPct\% interval} & \textbf{Verdict} \\
\midrule
\tabinput{../results/tables/b3}
\bottomrule
\end{tabularx}
\end{adjustwidth}
\end{table}

Every pass holds for $W$ as configured and with WL7's exclusions. Outside $W$ are the reserved but
unused part of each socket's receive memory and kernel memory outside the slab that no receive
queue is charged for. The co-tenants' growth of the merged cache is outside it too. In the
partial-ClientHello case nginx and Envoy left bytes queued, with $K_q$ of
\BThreePartialHelloNginxEpollCompetitorKq{} and \BThreePartialHelloEnvoyEpollCompetitorKq{} bytes
per pending connection on epoll and the same on io\_uring. So the two residuals of Section~\ref{sec:footprint} apply to their
passes. The server, in replay, left none queued. Tables~\ref{tab:b3parts} and~\ref{tab:b3partsserver}
give the parts of every $W$, the caches' growth and the field $f$.

The \CountLossBThree{} other cells in Holm are sslh-ev's, and each is a loss: the \CiPct\% interval
of $Q$ lies wholly below one. sslh-ev's counted footprint per pending connection was smaller than
the server's. In
the silent case $D$ is $\BThreeSilentSslhEvEpollDMedian$ bytes on epoll, with the interval
[$\BThreeSilentSslhEvEpollDLow$, $\BThreeSilentSslhEvEpollDHigh$], and
$\BThreeSilentSslhEvIouringDMedian$ on io\_uring, with
[$\BThreeSilentSslhEvIouringDLow$, $\BThreeSilentSslhEvIouringDHigh$]. On epoll the upper end of
$Q$'s interval prints as \BThreeSilentSslhEvEpollHigh{} at four places, below one before rounding.
At Holm's level that cell's interval, [\BThreeSilentSslhEvEpollHolmLow, \BThreeSilentSslhEvEpollHolmHigh],
contains one; the frozen definition of a loss reads the \CiPct\% interval.
In the partial-ClientHello case $D$ is $\BThreePartialHelloSslhEvEpollDMedian$ on epoll, with
[$\BThreePartialHelloSslhEvEpollDLow$, $\BThreePartialHelloSslhEvEpollDHigh$], and
$\BThreePartialHelloSslhEvIouringDMedian$ on io\_uring, with
[$\BThreePartialHelloSslhEvIouringDLow$, $\BThreePartialHelloSslhEvIouringDHigh$].

\Kbase{} is \KBase{} bytes per pending connection: the median of \ophold's $K_s$ over its valid
windows. Every system pays it alike, and it is most of the server's $W$, which is
\BThreeSilentNginxEpollServerW{} bytes in the silent case against nginx on epoll. There the server's
$U$ is \BThreeSilentNginxEpollServerU{} bytes, and its $K_s - \Kbase$ is
$\BThreeSilentNginxEpollServerKsNet$ bytes. $K_s$ falls below \Kbase{} in many arms,
every arm on epoll among them, so \Kbase{} is a shared reference, not a floor.

B3's \CountDescBThree{} descriptive cells are reported against the same bound and decide nothing.
Netty held less than the server in the partial-ClientHello case, with $Q$ of
\BThreePartialHelloNettyEpollMedian{} on epoll and \BThreePartialHelloNettyIouringMedian{} on
io\_uring and $D$ of $\BThreePartialHelloNettyEpollDMedian$ and $\BThreePartialHelloNettyIouringDMedian$
bytes. Every other descriptive cell with an interval lies above the bound. One cell,
cmux's partial-ClientHello cell on epoll, ended with \BThreePartialHelloCmuxEpollValid{} valid
sessions after its reruns, short of $R$, and so has no interval. Its
\BThreePartialHelloCmuxEpollInvalidWindows{} invalid windows each found the host's count of
TIME-WAIT sockets changed at a sample.

By the narrowing order, the B3 claim is narrowed to the competitors and cases it beats. These are
nginx, HAProxy and Envoy in both cases and hyper-util in the silent case, on both Linux backends.

\subsection{Mechanisms}\label{sec:res-mech}

Table~\ref{tab:mech} gives every cell. Of the \MHolm{} cells, \MPassed{} pass under both
computations, which disagree on \MDisagree{} cells.

\begin{table}[H]
\caption{The mechanism family. M1 is default / other detection mode in connections per second. M2
is in-process / relay in CPU per connection (WL6). M3 is the server's relay / the proxy in
connections per second.\label{tab:mech}}
\centering\scriptsize
\begin{tabularx}{\textwidth}{@{}l l l r l l X@{}}
\toprule
\textbf{Hyp.} & \textbf{Host} & \textbf{Backend, protocol, system} & \textbf{Valid} & \textbf{Median} &
\textbf{\CiPct\% interval} & \textbf{Verdict} \\
\midrule
\tabinput{../results/tables/mechanism}
\bottomrule
\end{tabularx}
\end{table}

\runin{M1.} M1 passes in \CountPassMOne{} of its \CountHolmMOne{} cells. Replay, the default on
every backend, served more connections per second than peek for HTTP/1.1 on all three backends
and for h2c on epoll and io\_uring. For h2c on IOCP a difference was not shown: the median is
\MOneWIocpHTwoCMedian, with the interval [\MOneWIocpHTwoCLow, \MOneWIocpHTwoCHigh]. On IOCP, peek
is a zero-byte receive followed by a synchronous peek. An undecided peek would switch to replay,
since Windows has no low-water mark. None did: the counters show \SwitchIocpHttpPeekSwitches{}
switches in \SwitchIocpHttpPeekAccepted{} connections of the valid peek windows for HTTP/1.1, and
\SwitchIocpHtwocPeekSwitches{} in \SwitchIocpHtwocPeekAccepted{} for h2c. So the IOCP cells compare
replay with that two-step peek, and measure the interface as much as the idea.

\runin{M2.} M2 passes in \CountPassMTwo{} of its \CountHolmMTwo{} cells, the two cells of HTTP/1.1.
In-process dispatch used \MTwoLEpollHttpMedian{} times the CPU per connection of relay on epoll and
\MTwoLIouringHttpMedian{} times on io\_uring, at the rates \RateMtwoEpoll{} and \RateMtwoIouring{}
exchanges per second. The two TLS cells were not run. In closed loop the back end that terminates
TLS limits the relay arm, so its relay windows broke the rule for back-end cores. Development runs
showed this before the freeze, and the author decided then not to run the two cells. The rate
sessions after the pilot entry ended the same way, so the frozen rule computes no rate for them.
Each cell enters Holm with \PIsOne, and a difference was not shown. M2's claim narrows to
HTTP/1.1.

\runin{M3.} M3 passes in \CountPassMThree{} of its \CountHolmMThree{} cells. The server's relay
served more connections per second than nginx, HAProxy, Envoy and caddy-l4, for TLS routed by its
server name and for HTTP/1.1. The medians range from \MThreeLEpollHttpNginxMedian{} against
nginx for HTTP/1.1 to \MThreeLEpollTlsSniEnvoyMedian{} against Envoy for TLS. The two cells of
sslh-ev are untested. Every one of their sslh-ev windows broke the error rule, as the defect that
the revision log records before the freeze predicts, and some also had failed connects. In
development the defect showed no failed connect and no listen overflow, and the cause of these is
not established. sslh-ev's listen
overflow counter rose by a total of \MThreeLEpollTlsSniSslhEvProxyListenOverflows{} and
\MThreeLEpollHttpSslhEvProxyListenOverflows{} in the two cells; no other proxy's rose. By the
narrowing order, M3 is narrowed to the four proxies it beats.

\runin{Secondary mechanism cells.} At a fixed load, replay's median time to first byte was below
peek's in every M1 cell (Table~\ref{tab:secmech}). The interval lies wholly below one in all but
the h2c cell on IOCP. Against the proxies the server's relay used less CPU per connection
in every cell, with medians from \SecMTtfbCpuMThreeLEpollHttpNginxCpuMedian{} against nginx to
\SecMTtfbCpuMThreeLEpollTlsSniEnvoyCpuMedian{} against Envoy. Its median time to first byte at a
fixed load was higher than nginx's, \SecMTtfbCpuSecMTtfbMThreeLEpollHttpNginxValueMedian{} times
for HTTP/1.1 and \SecMTtfbCpuSecMTtfbMThreeLEpollTlsSniNginxValueMedian{} for TLS, and lower than
the other three proxies'.

\subsection{The Secondary Cells}\label{sec:res-secondary}

Section~\ref{sec:secondary} lists the secondary cells; Tables~\ref{tab:sec} and~\ref{tab:secmech}
give each, and none decides anything. The results files list \CountNotRunList{} cells of any kind
short of their $R$, each with its logged reason: \CountNotRunNotRun{} were not run, and
\CountNotRunShort{} ran without reaching $R$. Section~\ref{sec:deviations} states why the secondary
cells ran again on L.

\runin{SSH.} Under churn the one-port arm served \SecSshSecSshCOneLEpollValueMedian{} and
\SecSshSecSshCOneLIouringValueMedian{} times the dedicated arm's connections per second on L, and
\SecSshSecSshCOneWIocpValueMedian{} on W. At a fixed load its median time to first byte was
\SecSshSecSshCThreeLEpollValueMedian, \SecSshSecSshCThreeLIouringValueMedian{} and
\SecSshSecSshCThreeWIocpValueMedian{} times dedicated mode's. The message order differs by mode, as
stated before the run: in one-port mode the server's identification line waits for the client's.
The \CountSshCtwoCells{} SSH cells of C2 could not run. The generator has no keep-alive form of SSH,
and the server's SSH handler closes after one line. Each cell ran its sessions and its reruns, and
all \CountSshCtwoWindows{} of its windows were invalid. On L the archived job \JobSlOne{} had shown
the same (Section~\ref{sec:dev-b1}).

\runin{The mixed-protocol cell.} On epoll the one-port arm served \SecMixedSecMixedCOneLEpollValueMedian{}
times the dedicated arm's connections per second, with the interval
[\SecMixedSecMixedCOneLEpollValueLow, \SecMixedSecMixedCOneLEpollValueHigh]. On IOCP it served
\SecMixedSecMixedCOneWIocpValueMedian{} times, with [\SecMixedSecMixedCOneWIocpValueLow,
\SecMixedSecMixedCOneWIocpValueHigh]. By design the background's silent connections live
differently in the two modes. In one-port mode the listener has no fallback, so the server closes
each at \Tdec{} and the holder opens it again at once. In dedicated mode they are held for the whole
window. The archived job \JobSlOne{} showed this in its counters. So the background's make-up is the
same in both modes, and its connections' lives are not.

On io\_uring the cell has no valid session. It ran \CountMixedIouringSessions{} sessions and
\CountMixedIouringWindows{} windows, of which \CountMixedIouringValid{} were valid, in both arms
alike. The probe of HTTP/1.1 timed out in \CountMixedIouringProbe{} windows, and the share of failed
exchanges exceeded its bound in \CountMixedIouringErrors. The churn starved beside its keep-alive
background in both modes, as \JobSlOne{} had shown, so this is not a cost of detection. We report it
as a measured limitation of the server's io\_uring backend. That backend's receives take buffers
that the kernel selects from a ring, which holds \RingBuffers{} free buffers per worker and gets a
new one each time the worker takes one. A receive that finds the ring empty fails and is posted
again. The cell's background keeps \NBg{} TLS and \NBg{} MQTT connections busy. How these might
starve the churn is a hypothesis, stated in Section~\ref{sec:disc-hyp}.

\runin{TLS variants.} With ALPN \texttt{h2}, the one-port arm served \SecTlsVariantsSecTlsVariantsCOneLEpollAlpnHTwoValueMedian,
\SecTlsVariantsSecTlsVariantsCOneLIouringAlpnHTwoValueMedian{} and
\SecTlsVariantsSecTlsVariantsCOneWIocpAlpnHTwoValueMedian{} times the dedicated arm's connections per
second on epoll, io\_uring and IOCP. TLS with session resumption was not run, for the reason given in
Section~\ref{sec:secondary}.

\runin{Two server cores.} With two cores the one-port arm served \SecTwoCoresSecTwoCoresCOneLEpollTwoCoresValueMedian{}
and \SecTwoCoresSecTwoCoresCOneLIouringTwoCoresValueMedian{} times the dedicated arm's connections
per second on epoll and io\_uring. The cell on IOCP was not run: the server serves one worker on
IOCP, and more would be a different threading model, outside this paper. A
\texttt{SO\_REUSEPORT} group of two listeners served \SecTwoCoresSecTwoCoresCOneLEpollReuseportValueMedian{}
times the shared listener's connections per second on epoll, with the interval
[\SecTwoCoresSecTwoCoresCOneLEpollReuseportValueLow, \SecTwoCoresSecTwoCoresCOneLEpollReuseportValueHigh].
On io\_uring it served \SecTwoCoresSecTwoCoresCOneLIouringReuseportValueMedian{} times, with
[\SecTwoCoresSecTwoCoresCOneLIouringReuseportValueLow, \SecTwoCoresSecTwoCoresCOneLIouringReuseportValueHigh].

\runin{The relay on io\_uring.} Against the four proxies with valid sessions, the server's relay on
io\_uring served from \SecRelayIouringSecRelayIouringLIouringHttpNginxValueMedian{} times nginx's
connections per second, for HTTP/1.1, to \SecRelayIouringSecRelayIouringLIouringTlsSniEnvoyValueMedian{}
times Envoy's, for TLS. sslh-ev's \CountSecSslh{} secondary cells ended with no valid session, as
its M3 cells did and as the follow-up of \JobSlOne{} expected. Its two cells of the time to first byte at a fixed
load ran closed-loop windows instead, since M3 gave them no rate.

\runin{IOCP's forms.} AcceptEx with a receive buffer served
\SecIocpFormsSecIocpFormsWIocpAcceptexBufferValueMedian{} times the default form's connections per
second, with [\SecIocpFormsSecIocpFormsWIocpAcceptexBufferValueLow,
\SecIocpFormsSecIocpFormsWIocpAcceptexBufferValueHigh]. The posted receive buffer served
\SecIocpFormsSecIocpFormsWIocpReceiveFormValueMedian{} times, with
[\SecIocpFormsSecIocpFormsWIocpReceiveFormValueLow, \SecIocpFormsSecIocpFormsWIocpReceiveFormValueHigh].
The posted form holds a buffer for a connection that has received no byte, so it breaks B2's memory
bound.

\runin{B3 in the other detection mode.} With the server in peek, its $W$ was
\SecBThreeOtherModeSecBThreeOtherModeEpollSilentValueMedian{} and
\SecBThreeOtherModeSecBThreeOtherModeIouringSilentValueMedian{} times its $W$ in replay in the silent
case, on epoll and io\_uring. In the partial-ClientHello case it was
\SecBThreeOtherModeSecBThreeOtherModeEpollPartialHelloValueMedian{} and
\SecBThreeOtherModeSecBThreeOtherModeIouringPartialHelloValueMedian{} times. Peek leaves the received
bytes queued in the kernel, where $K_q$ counts them.

\runin{Counts.} The operations and copies per connection, the system calls of the proxies and the
check of the server's system-call counters are in Section~\ref{sec:res-ops}.

\subsection{Operations, Copies and System Calls}\label{sec:res-ops}

These outputs are counts. None is timed and none decides anything. Two sources hold them. The
measured windows hold the server's counters over each window's process, and Section~\ref{sec:counters}
defines them. An untimed lab job, \JobSt, ran after the analysis on L's measured build. It ran the
frozen tool for these windows, after the freeze guard, and holds counts only
(Section~\ref{sec:dev-not-produced}). Its churn is open loop at \TraceRate{} exchanges per second for
\TraceSeconds~s after \TraceWarmSeconds~s, and its keep-alive uses \TraceKaConns{} connections for
\TraceKaSeconds~s. C1 and C3 share that churn, so their untimed windows differ only in name.

The copy counters count bytes, not buffers. ``Copied in user space'' counts each move of a payload
byte by the server's own code from one user-space place to another. The bytes that a system call
moves count as received, sent or peeked. Producing a byte, as a handler, OpenSSL or nghttp2 does,
is not a copy.

\runin{Operations and copies, in-process.} Appendix Tables~\ref{tab:ops-cost} and~\ref{tab:ops-m1}
give, per arm, the means over the valid windows of each window's counts, per connection, and per
request in C2. Table~\ref{tab:ops-cost} gives operations, bytes copied in user space and bytes
received and sent; Table~\ref{tab:ops-m1} adds bytes peeked and wakeups during detection.
Operations are the system calls and ring submissions the counters name. In replay the counter of
bytes copied in user space recorded none in any arm of the \WlFiveCostCells{} cost cells, and none
was peeked. The largest means are \WlFiveCostCopiedMax{} and \WlFiveCostPeekedMax{} bytes. One-port
mode's operations differed from dedicated mode's by at most \WlFiveCostOpsMaxDiffConn{} per
connection in C1 and C3 and \WlFiveCostOpsMaxDiffReq{} per request in C2. In M1, peek copied
\WlFiveMOneLEpollHttpPeekPeeked{} bytes per connection with \texttt{MSG\_PEEK}, which stay queued in
the kernel, on every backend.

\runin{Operations and copies, relay.} No measured window runs the relay in peek. So both detection
modes of the relay come from \JobSt, under one load, per relayed connection (Appendix
Table~\ref{tab:ops-relay}). Relay dispatch has no counterpart in dedicated mode, so peek is read
against replay. These counts depend on the load. At a low rate each step waits on its own, so the
loop's waits and the accepts that find no connection are more per connection. They are
therefore not those of M2's and M3's relay windows, which ran at other loads. The relay reads into its buffers and sends from them, so each relayed byte
counts once as received and once as sent. Its counter of bytes copied in user space recorded none
in either mode on either backend: the largest mean is \TraceRelayCopiedMax. In peek it peeked
\TraceRelayEpollHttpPeekPeeked{} bytes per relayed connection for HTTP/1.1 and
\TraceRelayEpollTlsSniPeekPeeked{} for TLS on epoll. Its operations per relayed connection for
HTTP/1.1 were \TraceRelayEpollHttpReplayOps{} in replay and \TraceRelayEpollHttpPeekOps{} in peek on
epoll, and \TraceRelayIouringHttpReplayOps{} and \TraceRelayIouringHttpPeekOps{} on io\_uring.

\runin{System calls of the proxies.} Table~\ref{tab:syscalls} gives each proxy's system calls per
connection, from \texttt{perf trace -s}, in front of the stub as in M3, with the server's relay
beside them. For HTTP/1.1 nginx made \TraceSysNginxHttpCalls{} calls per connection, HAProxy
\TraceSysHaproxyHttpCalls, Envoy \TraceSysEnvoyHttpCalls, caddy-l4 \TraceSysCaddyLFourHttpCalls{}
and sslh-ev \TraceSysSslhEvHttpCalls. The server's relay made \TraceSysRelayHttpCalls. For the calls
that the server's checks name, \texttt{perf stat}'s count stands beside that of \texttt{perf trace
-s}. In every row \texttt{perf trace -s} counted fewer of them, by \TraceSysShortfallMinPct\% to
\TraceSysShortfallMaxPct\%, and reported \TraceSysLost{} events lost.

\runin{The counter check.} The untimed windows cover the \TraceCostCells{} cost cells of L in both
modes. In them, \TraceCostChecksStatAgree{} of the \TraceCostChecks{} checks of the server's counters
equal \texttt{perf stat}'s count of the same calls. Appendix Table~\ref{tab:counter-check} gives them per
cell. That is the check the frozen text asks for, read before the freeze as a comparison with
\texttt{perf stat}. \texttt{perf trace -s} agreed in \TraceCostChecksTraceAgree, short of
\texttt{perf stat} by at most \TraceCostShortfallMaxPct\% of a call's count, with \TraceCostLost{}
events reported lost. In the relay's \TraceRelayRows{} rows, \TraceRelayChecksStatAgree{} of
\TraceRelayChecks{} checks equal \texttt{perf stat}'s count.

\begin{table}[H]
\caption{The proxies' system calls per connection on L, from \JobSt's untimed windows: open-loop
churn at \TraceRate{} exchanges per second, in M3's arrangement with the stub. ``All'' counts every call that
\texttt{perf trace -s} saw, over the connections the stub accepted. ``Checked'' counts only the
calls the server's checks name, by \texttt{perf trace -s} and by \texttt{perf stat}; the shortfall
is their difference over \texttt{perf stat}'s count.\label{tab:syscalls}}
\centering\scriptsize
\begin{tabularx}{\textwidth}{@{}l l r r r r r X@{}}
\toprule
\textbf{System} & \textbf{Protocol} & \textbf{Conns.} & \textbf{All} & \textbf{Checked, trace} &
\textbf{Checked, stat} & \textbf{Shortfall \%} & \textbf{Load ended ok} \\
\midrule
\tabinput{../results/tables/syscalls}
\bottomrule
\end{tabularx}
\end{table}

\subsection{The Competitors' Hard Cases}\label{sec:res-comp}

The competitors' table runs every hard case that a system's listener setup covers, at matched
timers and at the system's defaults, \CompReps{} replicates each. It is descriptive and decides
nothing, and no time from it is reported. Each row records whether the reply equals what the
server's frozen table expects. That field is neither a pass nor a failure of the system, whose own
specification may differ from the server's table.

The table holds \CompCaseRows{} rows from \CompCaseSystems{} systems, over \CompCaseVariants{}
variants of \CompCaseCases{} cases. Of the \CompCaseRowsObserved{} rows observed, the reply of
\CompCaseRowsAsExpected{} equals what the server's frozen table expects, and the reply of
\CompCaseRowsDiffer{} does not. \CompCaseRowsNotObserved{} rows were not observed, since the case
generator did not end. \CompCaseRowsNoDecision{} observed rows reached no decision within \TObsS~s
and are counted by their reply. The replicates agreed within every variant: \CompCaseUnitsMixed{} variants are mixed.
Table~\ref{tab:comp} gives the counts per system.

\begin{table}[H]
\caption{The competitors' hard cases, per system, on L. A variant is one case's form at one timer
setting with its replicates; it counts as equal or differing only when every replicate does.
``Equal'' means the reply equals what the server's frozen table expects.\label{tab:comp}}
\begin{adjustwidth}{-\extralength}{0cm}
\centering\scriptsize
\begin{tabularx}{\fulllength}{@{}l r r r r r r r r r X@{}}
\toprule
& \multicolumn{4}{c}{\textbf{Rows}} & & \multicolumn{2}{c}{\textbf{Matched timers}} &
\multicolumn{2}{c}{\textbf{Default timers}} & \\
\textbf{System} & \textbf{All} & \textbf{Equal} & \textbf{Differ} & \textbf{Not obs.} &
\textbf{Variants} & \textbf{Equal} & \textbf{Differ} & \textbf{Equal} & \textbf{Differ} &
\textbf{Cases not covered} \\
\midrule
\tabinput{../results/tables/competitors}
\bottomrule
\end{tabularx}
\end{adjustwidth}
\end{table}
